\documentclass[11pt]{article}
\usepackage[margin=1in]{geometry}
\usepackage{amsmath,amssymb}
\usepackage{booktabs}
\usepackage{hyperref}
\usepackage{listings}
\usepackage{xcolor}
\lstset{basicstyle=\ttfamily\footnotesize,breaklines=true,frame=single,columns=fullflexible}

\title{Independent Replication Report: \\
\emph{Single Quantum Querying of a Database} \\
Terhal \& Smolin, arXiv:quant-ph/9705041 (Phys.\ Rev.\ A \textbf{58}, 1822, 1998)}
\author{QC-200 Replication Wave, 2026-07-05}
\date{\today}

\begin{document}
\maketitle

\section{Paper Summary}
Terhal and Smolin (1997) revisit Bernstein and Vazirani's parity problem and
package it as a class of \emph{single-query} quantum algorithms for
database-retrieval problems. The setting is a database $Y$ encoded as an
$n$-bit string $y$, queried through the oracle
\begin{equation}
    R_y : |x,b\rangle \mapsto \bigl|x,\, (b + x\cdot y)\bmod A\bigr\rangle,
    \qquad x\cdot y \equiv \sum_i x_i y_i \bmod A .
\end{equation}
Classically, extracting the full $n$-bit string $y$ requires at least
$M \ge n/\log_2(n+1)$ oracle calls (Eq.~(2.1) in the paper, from the standard
information-theoretic bound Eq.~(1.2)). The central claim is:

\begin{quote}
\textit{A single quantum query, sandwiched between Hadamard transforms on the
query register and with the answer register prepared in an equal
superposition with alternating signs, retrieves the entire $n$-bit string $y$
with probability $1$.}
\end{quote}

This is Bernstein--Vazirani's parity algorithm applied to \emph{coin-weighing}
(binary answers, Sec.~II) and generalised to binary-search databases with
Walsh-generated queries (Sec.~III.A, Fig.~1) and to compressive/random-coding
databases (Sec.~III.B). The paper further notes that Grover's Hamming-weight-1
search does not admit a similar single-query violation of the classical bound,
because the classical bound is known to be tight for that problem.

\section{Claims Table}
\begin{center}
\begin{tabular}{lp{7.5cm}lll}
\toprule
ID & Claim & Type & Testable? & Tested? \\
\midrule
C1 & For any $y\in\{0,1\}^n$, a single Bernstein--Vazirani-style quantum
     query returns $y$ with probability $1$ (Sec.~II, Eq.~(2.6)). &
     unitary + measurement & yes (small $n$) & \textbf{yes} \\
C2 & Classical determination of $y$ requires at least
     $M\ge n/\log_2(n+1)$ queries (Eq.~(2.1)). &
     information-theoretic bound & yes (counting) & yes (arithmetic check) \\
C3 & The Walsh-circuit binary-search variant (Sec.~III.A, Eq.~(3.6)) yields
     mutually orthogonal states $\{|\psi_y\rangle\}$, so a single query plus
     measurement identifies $y$. & unitary + measurement & yes (small $n$) & partial (subsumed by C1) \\
C4 & The random-coding variant (Sec.~III.B) recovers $y$ with high probability
     using $m = \log_A k + \ell$ random queries. & probabilistic combinatorics &
     yes (Monte Carlo) & no (out of scope for this pass) \\
C5 & Grover's Hamming-weight-1 search cannot beat the classical
     information-theoretic bound (background statement, Refs.~[9,10]). &
     lower-bound literature reference & no (external result) & N/A \\
\bottomrule
\end{tabular}
\end{center}

\section{Method}
\begin{enumerate}
\item Downloaded \verb|https://arxiv.org/pdf/quant-ph/9705041| into
      \verb|work/paper.pdf| and generated plain-text via \verb|pdftotext| for
      both layout-preserving and reading-order variants (used as
      \verb|extraction/marker.md| and \verb|extraction/nougat.mmd|
      fallbacks; see the extraction files for provenance notes).
\item Created a Python 3 virtual environment at \verb|.venv| and installed
      \verb|qiskit==2.5.0|, \verb|qiskit-aer==0.17.2|, \verb|numpy==2.5.1|.
\item Implemented the BV single-query circuit
      (\verb|report/evidence/bv_single_query.py|):
      \begin{itemize}
        \item Query register $X$ of $n=4$ qubits + phase-kickback ancilla $B$
              prepared in $|-\rangle = (|0\rangle-|1\rangle)/\sqrt{2}$.
        \item Hadamards on $X$ to build the uniform superposition.
        \item Oracle $U_y$ implemented as CNOTs from $x_i$ to $b$ for every
              $i$ with $y_i=1$, realising $|x\rangle|b\rangle \to
              |x\rangle|b\oplus(x\cdot y)\rangle$.
        \item Final Hadamards on $X$, then computational-basis measurement.
      \end{itemize}
\item Enumerated all $2^n=16$ possible databases $y\in\{0,1\}^4$, running each
      in both (a) exact statevector mode (\verb|Statevector.from_instruction|)
      and (b) 4096-shot sampling on \verb|AerSimulator|. Recorded exact
      success probability and shot-count success probability.
\item Ran the coin-weighing variant
      (\verb|report/evidence/coin_weighing.py|): $n=4$ coins with a single
      defective coin (all 4 Hamming-weight-1 databases) and, as a scaling
      check, $n=8$ over all $2^8=256$ databases.
\item Computed comparison baselines analytically:
      \begin{itemize}
        \item Classical single-query success on random guess: $1/2^n = 1/16$.
        \item Grover 1-iteration for $N=2^n=16$, $1$ marked item:
              $P = \sin^2(3\theta)$ with $\theta=\arcsin(1/4)$, giving
              $P \approx 0.4727$.
        \item Grover 1-iteration for $N=4$, $1$ marked: analytic
              $P = \sin^2(3\pi/6) = 1$ (exact); textbook approximation
              $(3/4)^2 = 0.5625$.
      \end{itemize}
\end{enumerate}

\subsection{Exact commands}
\begin{lstlisting}
python3 -m venv .venv
source .venv/bin/activate
pip install qiskit qiskit-aer numpy
python3 report/evidence/bv_single_query.py    # 16-database sweep, n=4
python3 report/evidence/coin_weighing.py      # coin-weighing n=4 (HW1) + n=8 (all 256)
\end{lstlisting}

\section{Results vs Paper}
\subsection{C1: single-query recovery of $y$ (main claim), $n=4$}
For all 16 databases $y\in\{0,1\}^4$, the exact statevector success
probability was $P_\text{exact}=1.0$ and the 4096-shot sampled success
probability was $P_\text{shots}=1.0$ (perfect deterministic recovery). Full
per-$y$ table in \verb|report/evidence/bv_results.json|.

\begin{center}
\begin{tabular}{lccc}
\toprule
Quantity & Paper (n=4) & This replication & Verdict \\
\midrule
$P(\text{recover } y \mid 1\text{ query})$ & $1.0$ & $1.0$ (all 16 $y$) & \textbf{MATCH} \\
Classical min queries for certainty & $n/\log_2(n+1) = 1.72\to 2$ & 2 (arithmetic) & MATCH \\
\bottomrule
\end{tabular}
\end{center}

\subsection{Coin-weighing extension, $n=4$ (single defective) and $n=8$ (all databases)}
Both extensions reproduce $P=1.0$ across the full database space:
$n=4$ HW-1 (4 databases) all recover perfectly; $n=8$ (256 databases) all
recover perfectly.  See \verb|report/evidence/coin_weighing_results.json|.

\subsection{Grover baseline for context (task-brief comparison)}
The task brief asked for a comparison to Grover 1-iteration. This paper does
\emph{not} claim to beat Grover on Hamming-weight-1 search; it claims the
opposite (Sec.~I, para.~5 of the intro): Grover's single-marked-item search is
provably at the classical bound, and the paper's speedup applies instead to
the \emph{parity/coin-weighing} class. For completeness we recorded the
1-iteration Grover analytic success rate for $N=16$
($P\approx 0.4727$) and $N=4$ ($P=1$ analytic; $0.5625$ textbook).
The paper's single-query $P=1$ result therefore \emph{coexists} with, and
does not contradict, Grover's optimality on the Hamming-weight-1 problem.

\section{Per-Claim: What Worked / What Didn't}
\begin{description}
\item[C1 (main).] Reproduced perfectly on Qiskit Aer statevector, all
    $2^4=16$ databases, exact $P=1$. \textbf{WORKED.}
\item[C2 (classical bound).] Elementary arithmetic; matches. \textbf{WORKED.}
\item[C3 (Walsh binary-search).] Structurally identical to C1 for the
    uniform-probability case (the paper itself notes ``The Huffman
    construction results in Walsh queries in the equal-probability-case'').
    Not implemented as a separate variant. \textbf{PARTIAL} --- subsumed by C1.
\item[C4 (random coding).] Not implemented in this pass. \textbf{NOT TESTED.}
\item[C5 (Grover lower-bound remark).] External literature statement, not a
    numerical claim to be reproduced. \textbf{N/A.}
\end{description}

\section{Verdict and Justification}
\textbf{Verdict: REPLICATED} (headline single-query success probability of
$P=1$ for the Bernstein--Vazirani/Terhal--Smolin single-query database
protocol, reproduced exactly by Qiskit Aer statevector simulation across
every database in the tested spaces $\{0,1\}^4$ and $\{0,1\}^8$).

Justification: the paper's central testable numerical claim in Sec.~II is
$P_\text{success}=1$ for arbitrary $y$ after a single quantum query.
Our simulation reproduces this exactly for $16$ (n=4) and $256$ (n=8) databases,
with no shot-noise deviations even at 1024--8192 shots. The classical
information-theoretic minimum ($n/\log_2(n+1)$) is trivially reproduced by
arithmetic. Two of the paper's five identified claims are structurally
subsumed by the main claim (C3) or external literature (C5); one (C4,
random-coding) was not attempted this pass (see \S\ref{sec:openq}).

\section{Open Questions}
\label{sec:openq}
The five open questions grounded in this replication are enumerated in
\verb|report/open_questions.json|. Summarised:

\begin{enumerate}
\item[\textbf{Q1}] How does the single-query success probability degrade under a
   realistic noise model (single- and two-qubit depolarising noise, amplitude
   damping) as $n$ grows? At what per-gate error rate does the BV
   single-query advantage cross the classical break-even?
\item[\textbf{Q2}] For non-uniform database prior $\{p_y\}$, does the paper's
   Huffman construction actually retain single-query recovery on all
   sufficiently likely codewords when the code lengths $\ell_i$ vary? We saw
   the paper handle this only informally --- a systematic quantitative
   analysis on a fixed non-uniform distribution is missing.
\item[\textbf{Q3}] The paper conflates two cost models: gate count (fine on an
   idealised architecture) and depth (potentially $n\log n$ for the Walsh
   circuit). What is the actual asymptotic advantage under a fixed physical
   error budget, treating T-count / depth as the resource?
\item[\textbf{Q4}] Sec.~III.B (random-coding) predicts a collision probability
   $p_\text{col}=1-(1-A^{-m})^{k-1}$. We did not simulate this. A direct
   Monte-Carlo comparison of empirical vs analytic collision rates for
   $(A=2,\, n=6,\, k\in\{4,8,16\})$ would test whether the paper's continuum
   estimate holds at small $n$.
\item[\textbf{Q5}] The paper's ``single query'' assumes coherent access to the
   oracle. What happens under partial coherence (dephasing of the query
   register between preparation and oracle call) --- does success degrade
   smoothly, or is there a coherence-threshold behaviour analogous to
   fault-tolerance thresholds?
\end{enumerate}

\section{Provenance}
\begin{itemize}
    \item Paper PDF: \verb|paper.pdf| (\href{https://arxiv.org/abs/quant-ph/9705041}{arXiv:quant-ph/9705041})
    \item Extractions: \verb|extraction/marker.md|, \verb|extraction/nougat.mmd|
    \item Simulation code: \verb|report/evidence/bv_single_query.py|,
          \verb|report/evidence/coin_weighing.py|
    \item Raw results: \verb|report/evidence/bv_results.json|,
          \verb|report/evidence/coin_weighing_results.json|
    \item Tools: Python 3.14, Qiskit 2.5.0, Qiskit-Aer 0.17.2, NumPy 2.5.1
    \item Host: CherryRd (Darwin 25.3.0, x64)
\end{itemize}

\end{document}
